\chapter{Worked proofs and derivations}
We now move from cataloguing results to deriving selected statements.
Every example connects a small mathematical argument with a possible
analysis of electrophysiological recordings.

\section{Proving a set identity}
Fix arbitrary $x$. By definition of union,
$x\notin A\cup B$ means that $x\notin A$ and $x\notin B$,
which means $x\in A^c\cap B^c$.
Thus $(A\cup B)^c=A^c\cap B^c$.
Equality of sets was proved by equivalence of membership statements.

\section{Distinguishing relations from functions}
Let $U$ be a set of identified units and $T$ an observation interval.
Spike occurrences form a relation $R\subseteq U\times T$.
The relation is generally not a function $U\to T$ because
one unit may have multiple spike times. It does define
$F:U\to\mathcal P(T)$ by
$F(u)=\{t\in T:(u,t)\in R\}$.
The codomain $\mathcal P(T)$ is the set of subsets of $T$.

\section{Cauchy--Schwarz from an orthogonal projection}
For vectors $x,y$ in a real inner-product space, $y\ne0$,
\[
0\leq\left\|x-\frac{\langle x,y\rangle}{\|y\|^2}y\right\|^2
=\|x\|^2-\frac{\langle x,y\rangle^2}{\|y\|^2}.
\]
Rearrangement yields
$|\langle x,y\rangle|\leq\|x\|\|y\|$.
The complex case inserts conjugates as dictated by the
inner-product convention. On centered square-integrable random
variables, this yields
$|\Cov(X,Y)|^2\leq\Var(X)\Var(Y)$.
Thus $|\rho_{XY}|\leq1$ for positive variances, while a constant
variable gives undefined rather than zero correlation.

\section{Least squares as orthogonal projection}
For $X\in\R^{T\times p}$ with full column rank, minimize
$L(\beta)=\|y-X\beta\|_2^2$.
Differentiating gives
\[
\nabla L(\beta)=-2X^\top(y-X\beta),\qquad
X^\top X\widehat\beta=X^\top y.
\]
The fitted residual $r=y-X\widehat\beta$ is orthogonal to
every column of $X$. One can write
$\widehat\beta=(X^\top X)^{-1}X^\top y$ under the
rank assumption; numerically QR/SVD or stable linear solvers
are preferable to explicit inversion.

\section{PCA as a constrained maximization}
Center the rows of $X\in\R^{T\times N}$ with $T>1$ and
form $C=X^\top X/(T-1)$. The projected sample variance
along unit vector $v$ is $v^\top Cv$.
With a Lagrange multiplier $\lambda$,
\[
\mathcal L(v,\lambda)=v^\top Cv-\lambda(v^\top v-1),
\qquad Cv=\lambda v.
\]
By the spectral theorem, maximizing the Rayleigh quotient
selects a leading eigenvector; orthogonality constraints
yield subsequent components.
This proves PCA optimality relative to \emph{sample variance},
not physiological importance.

\section{Participation-ratio bounds}
For $r$ strictly positive covariance eigenvalues,
Cauchy--Schwarz gives
$(\sum_i\lambda_i)^2\leq r\sum_i\lambda_i^2$.
Because all cross products $\lambda_i\lambda_j$ are nonnegative,
$(\sum_i\lambda_i)^2\geq\sum_i\lambda_i^2$.
Therefore
\[
1\leq\frac{(\sum_i\lambda_i)^2}{\sum_i\lambda_i^2}\leq r.
\]
Equality on the right holds when positive eigenvalues are equal;
for $r>1$, the lower bound is approached if one eigenvalue
dominates and the remaining ones vanish relative to it.

\section{Why a Poisson mixture is overdispersed}
Let $G>0$, $\E G=1$, and $N\mid G\sim\Poisson(\mu G)$.
The conditional mean and variance both equal $\mu G$.
Total expectation yields $\E N=\mu$; total variance gives
\[
\Var(N)=\E[\mu G]+\Var(\mu G)
=\mu+\mu^2\Var(G).
\]
Thus $\Var(N)/\E N=1+\mu\Var(G)$.
This is a model-based result, not proof that actual response
variability comes from a particular neural gain mechanism
\cite{goris2014}.

\section{Coherence and spectral positive semidefiniteness}
Suppose the spectral matrix is
\[
S(f)=
\begin{pmatrix}
S_{xx}(f)&S_{xy}(f)\\
S_{xy}(f)^*&S_{yy}(f)
\end{pmatrix}\succeq0.
\]
The determinant must be nonnegative:
$S_{xx}S_{yy}-|S_{xy}|^2\geq0$.
If both auto-spectra are positive,
$0\leq |S_{xy}|^2/(S_{xx}S_{yy})\leq1$.
This proves a normalized spectral dependence bound, not
biophysical connectivity.

\section{Conditional point-process likelihood}
If $\lambda(t\mid\mathcal H_t)$ is a valid predictable
conditional intensity and $t_1,\ldots,t_n$ are observed events
in $(0,T]$, the point-process log-likelihood relative to a
specified reference measure is, up to reference terms,
\[
\ell=\sum_{i=1}^n\log\lambda(t_i\mid\mathcal H_{t_i})
-\int_0^T\lambda(s\mid\mathcal H_s)\,ds.
\]
The event term rewards model intensity at observed spikes;
the compensator penalizes excessive predicted spikes.
It follows by combining the small-bin event and no-event
probabilities under regularity conditions.

\section{Predictive influence is not direct physical causation}
Let $z_t$ be a latent driver and define observed series
$x_t=z_t+\epsilon_t$ and $y_t=z_{t-1}+\eta_t$.
Past $x$ can help predict $y$, yet this construction contains
no direct physical edge from $x$ to $y$.
Consequently Granger prediction depends on the chosen
variables and may reflect unobserved common inputs
\cite{bastos2016,cekic2018}.
% --- expanded edition (October 2026) ---
\section{The time-rescaling theorem for a conditional intensity}
\begin{theorem}
Let $N$ be a simple point process on $[0,T]$ with continuous, strictly
positive conditional intensity $\lambda(t\mid\mathcal H_t)$, and let
$\Lambda(t)=\int_0^t\lambda(s\mid\mathcal H_s)\,ds$. Then the rescaled
times $\Lambda(t_1)<\Lambda(t_2)<\cdots$ form a unit-rate Poisson process
on $[0,\Lambda(T)]$.
\end{theorem}
\begin{proof}[Proof via the likelihood]
By Result 76, the joint density of the first $n$ event times is
$\prod_{i=1}^n\lambda(t_i\mid\mathcal H_{t_i})\exp(-\Lambda(t_n))$.
Change variables to $z_i=\Lambda(t_i)-\Lambda(t_{i-1})$. Because
$\Lambda$ is differentiable with derivative $\lambda$, and because
$\Lambda(t_i)$ depends only on events before $t_i$, the Jacobian matrix
of the map $(t_1,\ldots,t_n)\mapsto(z_1,\ldots,z_n)$ is triangular with
diagonal entries $\lambda(t_i\mid\mathcal H_{t_i})$. The density of
$(z_1,\ldots,z_n)$ is therefore
$\exp(-\sum_iz_i)$ on the positive orthant: the $z_i$ are iid
exponential with mean one.
\end{proof}
The proof shows exactly where predictability is used: the triangular
Jacobian requires that the intensity at $t_i$ not depend on $t_i$ itself
or on later events. A model whose ``history'' term accidentally uses the
current bin's spike passes every goodness-of-fit test trivially.

\section{Gauss--Markov and whitening}
If $y=X\beta+\varepsilon$ with $\Cov\varepsilon=\sigma^2\Omega$, $\Omega\succ0$
known, set $\tilde y=\Omega^{-1/2}y$ and $\tilde X=\Omega^{-1/2}X$. The
transformed errors have covariance $\sigma^2I$, so by Result 31 the best
linear unbiased estimator is OLS on the transformed data:
\[
\hat\beta_{\rm GLS}=(X^\top\Omega^{-1}X)^{-1}X^\top\Omega^{-1}y.
\]
For AR(1) errors with coefficient $a$, $\Omega^{-1/2}$ is the
prewhitening filter $\tilde y_t=y_t-ay_{t-1}$ (with a corrected first
sample). Regression of a slowly varying LFP feature on behaviour with
OLS gives unbiased but inefficient coefficients and badly overconfident
standard errors; prewhitening fixes both if $a$ is well estimated.

\section{Positive semidefiniteness of empirical spectral matrices}
\begin{proposition}
For any data, the multitaper or Welch cross-spectral estimate
$\hat S(f)=\frac1K\sum_{k=1}^Kz_k(f)z_k(f)^\ast$, with $z_k(f)$ the vector
of tapered Fourier coefficients across channels, is Hermitian positive
semidefinite. Consequently the estimated coherence lies in $[0,1]$.
\end{proposition}
\begin{proof}
Each $z_kz_k^\ast$ is Hermitian with $a^\ast z_kz_k^\ast a=|z_k^\ast a|^2\ge0$; a
nonnegative combination preserves both properties. Apply the $2\times2$
determinant argument used for Result 69.
\end{proof}
With $K=1$ (a single taper and trial), $\hat S$ has rank one and the
estimated coherence equals one at every frequency for any pair of
signals. Coherence estimation requires averaging, and the number of
averaged terms sets the bias floor $\approx1/K$.

\section{An identity behind trial-averaged evoked responses}
Let $x_r(t)=s(t)+n_r(t)$ for trials $r=1,\ldots,R$ with zero-mean noise
of variance $\sigma^2(t)$, independent across trials. The trial average
$\bar x$ has variance $\sigma^2/R$, and
\[
\E\Bigl[\frac1{R-1}\sum_r(x_r-\bar x)^2\Bigr]=\sigma^2,\qquad
\E[\bar x^2]=s^2+\frac{\sigma^2}R.
\]
Hence the power of the average overestimates the power of the evoked
signal by $\sigma^2/R$, and the unbiased estimate is
$\bar x^2-\hat\sigma^2/R$. Comparing evoked power between conditions with
different trial counts without this correction biases the comparison
towards the condition with fewer trials.
